\uuid{1jGQ}
\titre{Variables continues}
\auteur{worms}
\organisation{exo7}
\licence{CC BY-NC-SA 4.0}
\datecreate{2017-08-28}
\niveau{L2}
\module{Probabilité et statistique}
\chapitre{Probabilité continue}
\sousChapitre{Densité de probabilité}
\difficulte{}
\theme{Variables continues}
\exo7qcmid{worms.worms.039}

\contenu{
\begin{question}
\qclassification{261.01}
Soit $X$ le temps de trajet quotidien de Katrin, en heures, variable aléatoire de densité définie par
$$ f_X(x) = x\, e^{-x} \mathbb{I}_{[0,+\infty[}(x) $$
Laquelle des affirmations suivantes est vraie ?
\begin{answers}
\bad{$\mathbb{E}(X)=1$ et $\mathbb{P}(X>1)\simeq 63.2\%$.}
\bad{$\mathbb{E}(X)=2$ et $\mathbb{P}(X>1)\simeq 26.4\%$.}
\good{$\mathbb{P}(X>1)\simeq 73.6\%$ et on ne peut pas facilement déterminer la médiane de $X$.}
\bad{$\mathbb{E}(X)$ vaut $+\infty$ (c'est-à-dire $X$ n'admet pas d'espérance) et $\mathbb{P}(X>1)\simeq 26.4\%$.}
\end{answers}
\begin{explanations}
$\mathbb{E}(X)$ vaut effectivement $2$, mais $F(x)=1-(1+x)e^{-x}$ donc $\mathbb{P}(X>1)=\bar F(1) \simeq 73.6\%$. Quant à la médiane, $F$ est bien strictement croissante (le vérifier), mais l'équation $F(x)=1/2$ n'a pas de résolution explicite, il faut recourir à un schéma numérique pour la résoudre.
\end{explanations}
\end{question}
}
